%!TEX program = uplatex
\RequirePackage{plautopatch}
\documentclass[uplatex,dvipdfmx]{jlreq}
\usepackage{amsmath,amssymb,amscd,amsthm}

\pagestyle{empty}

\begin{document}

%======================================================================
% 講演１：山田光太郎「極小曲面論入門---Weierstrass表現公式からわかること」
%======================================================================
\begin{center}
{\Large\bfseries 極小曲面論入門---Weierstrass表現公式からわかること}

\hspace*{\fill}{\large\bfseries 山田 光太郎（九大・数理）}
\end{center}

\bigskip

3次元ユークリッド空間内の極小曲面（平均曲率が恒等的に $0$ であるような曲
面）は、古くから幾何学の重要な研究対象であった。
Weierstrass 表現公式は、この極小曲面を、曲面上の適当な複素構造に関する複
素解析的なデータで表す公式である。
極小曲面が解析関数で表されるということがわかればそれで十分、という考え方
もあるだろうが、一方、表現公式によって極小曲面の性質、とくに大域的な性質
を精密に調べられる、と思うこともできる。

本講演では極小曲面論入門として、Weierstrass 表現公式と、それから導かれる
極小曲面のいくつかの性質---
とくに完備極小曲面のエンドの形状と、フラックス問題---について解説する。
さらに、Weierstrass と類似の表現公式によって、解析関数で記述される曲面に
ついて言及したい。

\newpage

%======================================================================
% 講演２：小磯深幸「Plateau 問題について」
%======================================================================
\begin{center}
{\Large\bfseries Plateau 問題について}

\hspace*{\fill}{\large\bfseries 小磯 深幸（京教大・教育）}
\end{center}

\bigskip

3 次元ユークリッド空間 ${\mathbf R}^3$ 内の曲面 $M$ が極小曲面であるとは、
$M$ 上至るところ平均曲率が $0$ であるときをいう。
これは、$M$ の各点の近傍 $U$ について、
$U$ が境界を保つ任意の変分に対して面積汎関数の臨界点になる
ということと同値である。
それゆえ、しばしば極小曲面は石鹸膜の幾何学的モデルを与えると言われる。
石鹸膜は薄いので、その重力を無視して表面張力の作用だけを考えることにすれば、
面積の極小値を与えるような形をとると見なせるからである。

ところで、あらかじめ与えられた枠（たとえば1つのJordan閉曲線）に対して
それを境界とする極小曲面を求めるという問題は、
今日、Plateau 問題と呼ばれている。
これは、19世紀後半に石鹸膜等の実験により
この問題をさかんに研究したベルギーの物理学者 J.~Plateau にちなむものである。
Plateau 問題の解の存在が任意の Jordan 閉曲線に対して証明されたのは
1930 年のことで、J.~Douglas と T.~Radó によって独立にであった。
Douglas はこの研究により、第 1 回のフィールズ賞を受賞している。

本講演では、極小曲面及び関連する幾何学的変分問題の研究についての
歴史的なこと、
Plateau 問題の解決と解の性質、関連する最近の話題等について、
二回にわたって紹介する。
それぞれの回の内容としては次のものを予定している。

\medskip

1. 極小曲面論、主として Plateau 問題の古典解に関する入門的な講義

\smallskip

\quad ○ 歴史的なこと
--- 変分法の始まり、極小曲面論の発展（20世紀初頭まで）---

\quad ○ J.~Douglas と T.~Radó による Plateau 問題の解決
--- 円板型面積最小解の存在 ---

\quad ○ Plateau 問題の古典解の基本的な性質
--- 円板型面積最小解の regularity ---

\medskip

2. Plateau 問題の解に関する基本的かつ重要な研究、及び、関連する最近の話題

\smallskip

\quad ○ Plateau 問題の解の一意性・有限性

\quad ○ 極小曲面の安定性（面積極小性）

\quad ○ 関連する最近の話題について

\bigskip

さて、極小曲面の研究の歴史は、J.~L.~Lagrange が ``Essai d'une nouvelle
méthode pour déterminer les maxima et les minima des formules
intégrales indéfinies'' (1762) の中で、極値問題の例として
「${\mathbf R}^3$ 内に与えられた１つの閉曲線に対し、これを境界にもつ曲面全体の中
で面積最小のものをみつけたい」とし、面積最小曲面が $z = z(x, y)$ とグラフの
形で表されているときにそれが満たすべき方程式を導いたことに始まる。
同じ頃、すなわち18世紀後半から19世紀初頭にかけて研究が始まった曲面に関する変分問題
としては、平均曲率一定曲面（体積を保つ変分に対する表面積の臨界点）、
capillary 曲面（体積を保つ変分に対する表面積プラス重力エネルギーの臨界点）、
Willmore 曲面（平均曲率の2乗の積分の臨界点）等があるが、
いずれも物理現象にその動機を得ており、今日でも盛んに研究されている。

Lagrange の後19世紀後半までに、J.~B.~M.~C.~Meusnier、G.~Monge、A.~Legendre、
S.~F.~Lacroix、A.~M.~Ampère、H.~F.~Scherk、O.~Bonnet、J.~A.~Serret、
B.~Riemann、K.~Weierstrass、A.~Enneper、H.~A.~Schwarz 等により、
極小曲面に関する多くの基本的かつ重要な研究が行われた。
19世紀においては、極小曲面の研究の主な手法は函数論（複素解析）であった。
極小曲面を等温座標によって表示したときに各座標関数が調和関数になること
から、函数論が有効であったのである。
19世紀における極小曲面の研究の主な目的は、曲面の具体的な方程式を求めることで
あったが、
20世紀の数学の発展、とりわけ解析学の発展は、
極小曲面の研究に多様な視点と豊かな実りをもたらした。
1960年代後半以降には、解の一意性、解の正則性や微分可能性、
安定性、極小曲面の成す空間の構造などが、関数解析学、微分方程式論、
幾何学的測度論等の発展と呼応して新たな問題として現れた。

本講演では、Plateau 問題に関連する未解決問題にも言及すると共に、
今後の研究の方向性をも探りたい。

\newpage

%======================================================================
% 講演３：梅原雅顕「極小曲面のOsserman 不等式をめぐって」
%======================================================================
\begin{center}
{\Large\bfseries 極小曲面のOsserman 不等式をめぐって}

\hspace*{\fill}{\large\bfseries 梅原 雅顕（広島大・理）}
\end{center}

\bigskip

２次元リーマン多様体 $M^2$ についてその Gauss 曲率 $K$ の総和
$\displaystyle\int_{M^2}K\,dA$
を全曲率という。通常、コンパクトな多様体 $M^2$ については、Gauss--Bonnet の定理
\[
\frac{1}{2\pi}\int_{M^2}K\,dA=\chi(M)
\]
が成り立つ。さらに一般に $M^2$ がコンパクトとは
限らないが、完備な場合には、
「有限全曲率である」という仮定のもとで
Cohn-Vossen の不等式
\[
\frac{1}{2\pi}\int_{M^2}K\,dA\le \chi(M)
\]
が成り立つ。

ところで、
$\mathbf R^3$ における有限全曲率をもつ完備な極小曲面
を考えた場合、通常の Cohn-Vossen の不等式より
遙かに強い以下の Osserman の不等式
\[
\frac{1}{2\pi}\int_{M^2}K\,dA\le \chi(M)-\mbox{（エンドの数）}
\]
が成り立つ。
極小曲面の場合には、平均曲率が０であることから、
各点における
２つの主曲率 $\lambda_1,\lambda_2$ の符号は異なり、
その積である Gauss 曲率 $K=\lambda_1\lambda_2$ は
非正であり、もしも、この値が有限であるとすると、
実は、その値は $-4\pi$ の自然数倍であること
に注意する。
また、Osserman の不等式の等号は、曲面のすべての
エンドが自己交叉をしないことと同値であることが、
Jorge と Meeks によって示されている。
本講演では、以下の事柄について解説したい。
\begin{itemize}
\item[$\S1$] 不等式の証明と等号条件の証明の概略。
\item[$\S2$] $\mathbf R^n$ の極小曲面に対する一般化である
Chern--Osserman 不等式と等号条件。
\item[$\S3$] 不等式の解釈。
\item[$\S4$] 極小曲面以外の曲面に対する不等式の類似の紹介。
\item[$\S5$] ある種の特異点を許容する曲面の族について。
\end{itemize}
$H^3$ の平坦な曲面や、３次元 Minkowski 空間
$(\mathbf R^3,(-+++))$ の極大曲面については、完備な例が
自明なものしかないことが知られている。
実は、ある種の特異点を許容するこのような曲面の一般化
を考えると、完備な例が豊富に存在し、さらに
そのような曲面のカテゴリーにおいて
Osserman の不等式の類似が成立する。
このことに着目し、最後に、願望もこめて、
今後「特異点つきの曲面論」が、微分幾何の
稔り豊かな研究対象となりうることに言及したい。

\newpage

%======================================================================
% 講演４：宮岡礼子「極小曲面のガウス写像」
%======================================================================
\begin{center}
{\Large\bfseries 極小曲面のガウス写像}

\hspace*{\fill}{\large\bfseries 宮岡 礼子（上智大・理工）}
\end{center}

\bigskip

ここでは極小曲面の曲り方と、ガウス写像に関する話題を、
グラフィックを交えて紹介します。

全平面上のグラフ $(x,y,f(x,y))$ として与えられる
極小曲面が平面しかないことは、ベルンシュタインの定理として
よく知られています。
極小曲面の座標関数が調和関数であることと、
リュービルの定理からの帰結です。

極小曲面のガウス曲率は正にはなれない、つまり、極小曲面は
いたるところ鞍型の曲面です。これをグラフとしてかこうとすれば、
遠くへ行くほど平らにせざるを得ず、実際このときガウス曲率の
絶対値は原点からの距離の逆数で上から押さえられます。

曲面の各点に単位法ベクトルを対応させるのがガウス写像ですが、
グラフの場合、単位法ベクトルは常に赤道より上の半球面に値をとります。
Osserman は、任意の完備極小曲面のガウス写像について、
その除外値集合のキャパシティが大きいと、曲面は平面になることを示しました。
藤本坦孝氏は1988年、平面以外の完備極小曲面に対して、
この除外値集合は高々４点からなるという最良の結果を得ました。
リュービルの定理の精密化である Nevanlinna 理論をさらに発展させて
得られる、幾何学と複素関数論にまたがる重要な結果です。

これに関連する未解決問題と、そのアプローチについても話せればと考えています。

\end{document}
